\section{Related work}
\label{sec:related}

\paragraph{Adaptive loss weighting in PINNs.}
The fragility of the composite-loss objective in the original PINN
formulation \citep{raissi2019} has motivated a line of methods that
adapt the residual / boundary / data weights during training.
Self-adaptive PINN (SA-PINN) of \citet{mcclenny2020} attaches a
trainable per-collocation-point weight that is updated by ascent on
the residual, yielding an effective focusing of the network on
hard-to-fit regions; their ``points-with-weights'' formulation is
essentially a soft hard-example-mining for the residual. ReLoBRaLo
\citep{bischof2021} re-balances loss-term weights from running
ratios of the absolute loss values, smoothed by an exponential moving
average. Causal-PINN \citep{wang2024} weights the residual by a
time-causal mask so the network does not collapse onto late-time
dynamics before fitting the initial-time region. These all produce
\emph{one} solution for \emph{one} (adaptively-found) weighting and
would need to be retrained at every new weighting the practitioner
wishes to inspect.

\paragraph{Operator learning.}
A parallel line of work has built neural operators that map a
function-valued input (boundary condition, forcing term, initial
condition) to a function-valued solution. The Fourier Neural Operator
(FNO) \citep{li2020fno} parameterises the operator in spectral space:
linear layers on the lowest Fourier modes, plus a residual
$1\!\times\!1$ convolution. DeepONet \citep{lu2021deeponet} factorises
the operator into a branch network (encoding the input function at
sensor points) and a trunk network (encoding the spatial coordinate),
recombined by an inner product. Both methods amortise across a family
of problems but require a precomputed dataset of (input, solution)
pairs from a classical solver---they are supervised learners in
function space. The PINN residual itself plays no role in their
training. PI-DeepONet \citep{wang2021pi-deeponet} re-introduces the
residual but stays inside the operator-learning framing, treating the
input function as the parameter rather than a scalar coefficient.
Our LC-PINN sits between these two camps: it amortises like
operator learning but uses the residual loss like a PINN.

\paragraph{Loss-conditional networks.}
\citet{dosovitskiy2020} introduced the loss-conditional networks idea
in image generation: a single generator, conditioned on the loss
weights of a multi-objective composite reconstruction loss, produces
images that interpolate the tradeoff surface (sharpness vs.\ perceptual
fidelity vs.\ adversarial realism). The network is trained by sampling
the weight vector from a fixed prior at every step and applying the
freshly-sampled weights to the per-instance reconstruction loss.
LC-PINN takes the same conditional-on-weights idea and applies it to
the PDE-residual setting; the formal $\lambda$-invariance result
(\Cref{sec:theorem}) is, to our knowledge, the first analysis of the
optimum-set structure of this construction in the residual-loss
case.

\paragraph{Hypernetworks and meta-learning for PINNs.}
A separate amortisation strategy uses hypernetworks
\citep{ha2017hypernetworks} or MAML-style meta-learning
\citep{finn2017maml} to produce per-instance PINN weights. PI-MAML
\citep{liu2022pi-maml} demonstrates the meta-learning angle on
parametric PDEs. These methods amortise at the level of network
weights rather than a network input. They typically require a training
distribution of full PDE \emph{tasks} and a per-task inner-loop
adaptation step at test time; LC-PINN trains a single network with no
inner loop and produces predictions for every $\lambda$ in one forward
pass.

\paragraph{Extrapolation in $\lambda$ vs.\ in physical parameters.}
In the loss-weight mode the network's role is closest to
\citet{dosovitskiy2020}; in the parametric-coefficient mode it is
closest to FNO / DeepONet but parameterised over a low-dimensional
scalar input rather than a function-valued one. The two modes share
an architecture; the choice of $\lambda$ is task-driven.
